% GENERATED FILE. Edit canonical lessons, then run npm run generate:books.
% canonical-source-hash: fnv1a64:044329751e2b4865
% canonical-lessons: RU-C517-khor, RU-C517-orkestr, RU-C517-dirizhyor, RU-C517-dirizhirovat, RU-C517-kompozitor

\chapter{Choir and Orchestra}
\label{ch:ru-choir-and-orchestra}
\hlchaptermodality{517}

\begin{chapteropening}
\textbf{By the end of this chapter:} \emph{I can talk about choirs, orchestras and composers.}

Complete the last lesson of chapter 517: I can talk about choirs, orchestras and composers.
\end{chapteropening}

\section[\texorpdfstring{\ru{хор} (khor)}{хор (khor)}]{\ru{хор} (khor) — a choir}

\label{lesson:RU-C517-khor}



\begin{quote}
Before the new one: say the Russian for to hum, then the Russian for to sing out of tune.
\end{quote}

\subsection*{You'll want to know: \ru{хор}}
\textbf{\ru{хор}} (\emph{khor}) — \textquotedblleft{}a choir\textquotedblright{}.

\begin{quote}
\ru{В} \ru{школе} \ru{я} \ru{пел} \ru{в} \ru{хоре}.
\end{quote}

\begin{quote}
\emph{V shkóle ya pel v khóre.}
\end{quote}

\begin{quote}
At school I sang in the choir.
\end{quote}

Masculine.

\begin{cousinweb}
\emph{\ru{хор}} comes from Greek \emph{khorós} 'a company of dancers or singers' (EtymDB).
\end{cousinweb}

\begin{grammarlens}[title={What you've built}]
One more word for telling, explaining and describing at B1.
\end{grammarlens}

\subsection*{Your turn}
\begin{itemize}
\raggedright
  \item \emph{Say it:} \emph{khor}
  \item \emph{Say it:} \emph{khor}, once more
  \item \emph{Say it:} a sentence of your own with \emph{khor}
  \item \emph{Recall:} say the Russian for to hum, then the Russian for to sing out of tune, then say \emph{khor} again
\end{itemize}

\begin{tcolorbox}[breakable,colback=teal!4,colframe=teal!35!black,title={Before you move on}]
What does \ru{хор} mean? (\textquotedblleft{}A choir\textquotedblright{}.) Say it once more.
\end{tcolorbox}
\section[\texorpdfstring{\ru{оркестр} (orkéstr)}{оркестр (orkéstr)}]{\ru{оркестр} (orkéstr) — an orchestra}

\label{lesson:RU-C517-orkestr}



\begin{quote}
Before the new one: say the Russian for to sing out of tune, then the Russian for a choir.
\end{quote}

\subsection*{You'll want to know: \ru{оркестр}}
\textbf{\ru{оркестр}} (\emph{orkéstr}) — \textquotedblleft{}an orchestra\textquotedblright{}.

\begin{quote}
\ru{Оркестр} \ru{играл} \ru{два} \ru{часа} \ru{без} \ru{перерыва}.
\end{quote}

\begin{quote}
\emph{Orkéstr igrál dva chasá bez pererýva.}
\end{quote}

\begin{quote}
The orchestra played for two hours without a break.
\end{quote}

Masculine.

\begin{cousinweb}
\emph{\ru{оркестр}} comes from Greek \emph{orkh\'{\={e}}stra}, the dancing floor of a theatre, built on \emph{orkhéomai} 'to dance' (EtymDB).
\end{cousinweb}

\begin{grammarlens}[title={What you've built}]
One more word for telling, explaining and describing at B1.
\end{grammarlens}

\subsection*{Your turn}
\begin{itemize}
\raggedright
  \item \emph{Say it:} \emph{orkéstr}
  \item \emph{Say it:} \emph{orkéstr}, once more
  \item \emph{Say it:} a sentence of your own with \emph{orkéstr}
  \item \emph{Recall:} say the Russian for to sing out of tune, then the Russian for a choir, then say \emph{orkéstr} again
\end{itemize}

\begin{tcolorbox}[breakable,colback=teal!4,colframe=teal!35!black,title={Before you move on}]
What does \ru{оркестр} mean? (\textquotedblleft{}An orchestra\textquotedblright{}.) Say it once more.
\end{tcolorbox}
\section[\texorpdfstring{\ru{дирижёр} (dirizhyór)}{дирижёр (dirizhyór)}]{\ru{дирижёр} (dirizhyór) — a conductor (of an orchestra)}

\label{lesson:RU-C517-dirizhyor}



\begin{quote}
Before the new one: say the Russian for a choir, then the Russian for an orchestra.
\end{quote}

\subsection*{You'll want to know: \ru{дирижёр}}
\textbf{\ru{дирижёр}} (\emph{dirizhyór}) — \textquotedblleft{}a conductor (of an orchestra)\textquotedblright{}.

\begin{quote}
\ru{Дирижёр} \ru{поднял} \ru{палочку}, \ru{и} \ru{в} \ru{зале} \ru{стало} \ru{тихо}.
\end{quote}

\begin{quote}
\emph{Dirizhyór pódnyal pálochku, i v zále stálo tíkho.}
\end{quote}

\begin{quote}
The conductor raised his baton and the hall went quiet.
\end{quote}

Masculine.

\begin{cousinweb}
\emph{\ru{дирижёр}} is usually explained as built on French \emph{diriger} 'to direct'; EtymDB has no entry.
\end{cousinweb}

\begin{grammarlens}[title={What you've built}]
One more word for telling, explaining and describing at B1.
\end{grammarlens}

\subsection*{Your turn}
\begin{itemize}
\raggedright
  \item \emph{Say it:} \emph{dirizhyór}
  \item \emph{Say it:} \emph{dirizhyór}, once more
  \item \emph{Say it:} a sentence of your own with \emph{dirizhyór}
  \item \emph{Recall:} say the Russian for a choir, then the Russian for an orchestra, then say \emph{dirizhyór} again
\end{itemize}

\begin{tcolorbox}[breakable,colback=teal!4,colframe=teal!35!black,title={Before you move on}]
What does \ru{дирижёр} mean? (\textquotedblleft{}A conductor (of an orchestra)\textquotedblright{}.) Say it once more.
\end{tcolorbox}
\section[\texorpdfstring{\ru{дирижировать} (dirizhírovat)}{дирижировать (dirizhírovat)}]{\ru{дирижировать} (dirizhírovat) — to conduct (music)}

\label{lesson:RU-C517-dirizhirovat}



\begin{quote}
Before the new one: say the Russian for an orchestra, then the Russian for a conductor.
\end{quote}

\subsection*{You'll want to know: \ru{дирижировать}}
\textbf{\ru{дирижировать}} (\emph{dirizhírovat}) — \textquotedblleft{}to conduct (music)\textquotedblright{}.

\begin{quote}
\ru{Сегодня} \ru{оркестром} \ru{дирижирует} \ru{молодая} \ru{женщина}.
\end{quote}

\begin{quote}
\emph{Segódnya orkéstrom dirizhíruyet molodáya zhénshchina.}
\end{quote}

\begin{quote}
Today a young woman is conducting the orchestra.
\end{quote}

Imperfective; the present is \ru{дирижирую}. Takes the instrumental: \emph{\ru{дирижировать} \ru{оркестром}}.

\begin{cousinweb}
\emph{\ru{дирижировать}} is usually traced to French \emph{diriger} 'to direct', with \emph{-\ru{ировать}}; EtymDB lists the verb without a source.
\end{cousinweb}

\begin{grammarlens}[title={What you've built}]
One more word for telling, explaining and describing at B1.
\end{grammarlens}

\subsection*{Your turn}
\begin{itemize}
\raggedright
  \item \emph{Say it:} \emph{dirizhírovat}
  \item \emph{Say it:} \emph{dirizhírovat}, once more
  \item \emph{Say it:} a sentence of your own with \emph{dirizhírovat}
  \item \emph{Recall:} say the Russian for an orchestra, then the Russian for a conductor, then say \emph{dirizhírovat} again
\end{itemize}

\begin{tcolorbox}[breakable,colback=teal!4,colframe=teal!35!black,title={Before you move on}]
What does \ru{дирижировать} mean? (\textquotedblleft{}To conduct (music)\textquotedblright{}.) Say it once more.
\end{tcolorbox}
\section[\texorpdfstring{\ru{композитор} (kompozítor)}{композитор (kompozítor)}]{\ru{композитор} (kompozítor) — a composer}

\label{lesson:RU-C517-kompozitor}



\begin{quote}
Before the new one: say the Russian for a conductor, then the Russian for to conduct.
\end{quote}

\subsection*{You'll want to know: \ru{композитор}}
\textbf{\ru{композитор}} (\emph{kompozítor}) — \textquotedblleft{}a composer\textquotedblright{}.

\begin{quote}
\ru{Этот} \ru{композитор} \ru{пишет} \ru{музыку} \ru{для} \ru{кино}.
\end{quote}

\begin{quote}
\emph{Étot kompozítor píshet múzyku dlya kinó.}
\end{quote}

\begin{quote}
This composer writes film music.
\end{quote}

Masculine.

\begin{cousinweb}
\emph{\ru{композитор}} was borrowed from Latin \emph{compositor} 'an arranger' (EtymDB), from \emph{compōnere} 'to put together' (Lewis \& Short).
\end{cousinweb}

\begin{grammarlens}[title={What you've built}]
One more word for telling, explaining and describing at B1.
\end{grammarlens}

\subsection*{Your turn}
\begin{itemize}
\raggedright
  \item \emph{Say it:} \emph{kompozítor}
  \item \emph{Say it:} \emph{kompozítor}, once more
  \item \emph{Say it:} a sentence of your own with \emph{kompozítor}
  \item \emph{Recall:} say the Russian for a conductor, then the Russian for to conduct, then say \emph{kompozítor} again
\end{itemize}

\begin{tcolorbox}[breakable,colback=teal!4,colframe=teal!35!black,title={Before you move on}]
What does \ru{композитор} mean? (\textquotedblleft{}A composer\textquotedblright{}.) Say it once more.
\end{tcolorbox}
